\section{Dynamical zeta functions of fiber bundles over \texorpdfstring{$S^{1}$}{}}\label{section3}

\subsection{Cohomological expression}
We define the dynamical zeta function of $M$ by
\begin{equation*}
    \zeta(M,\Fh_{M},\phi^{t};s):=\prod_{\gamma}(1-N(\gamma)^{-s})^{-1}
\end{equation*}
where $\gamma$ runs over periodic orbits of the $\R$-action $\phi^{t}$.
Denote simply by $\zeta(M;s)$ the dynamical zeta function of $(M,\Fh_{M},\phi^{t})$.
C. Deninger suggested that the dynamical zeta function can be described in terms of leafwise cohomology groups and infinitesimal generator and is an analogue of the Hasse-Weil zeta function.
We show that $\zeta(M;s)$ has the leafwise cohomological expression for $(M,\Fh_{M},\phi^{t})$ as C. Deninger conjectured.

\begin{theorem}
The dynamical zeta function $\zeta(M;s)$ is described in terms of leafwise cohomology groups and infinitesimal generator 
    \begin{equation*}
        \zeta(M;s)=\prod_{i}^{d}\det_{\infty}(s\cdot\mbox{id}-\Theta|H^{i}_{\Fh}(M,\C))^{(-1)^{i+1}}
    \end{equation*}
where $\det_{\infty}$ is the regularized determinant and $H^{i}_{\Fh}(M,\C)$ is the complexification of $H^{i}_{\Fh}(M)$.
\end{theorem}


We first see that the leafwise cohomological expression is defined for $(M,\Fh_{M},\phi^{t})$. 
Recall that regularized determinants $\det_{\infty}(\Theta|V)$ is defined if the following condition (1) and (2) hold:

(1) $V$ is a complex vector space of countable dimension.
An operator $\Theta$ has finitely many eigenvalues whose eigen-space is of finite dimension. Then $V$ is the direct sum of eigen-spaces.

(2) Under condition (1), let $\Sp(\Theta|V)$ be the set of eigenvalues of $\Theta$ with multiplicities.
We assume that the Dirichlet series
\begin{equation*}
    \sum_{\alpha\ne0\in\Sp(\Theta|V)}\frac{1}{\alpha^{s}}\mbox{ with }\alpha^{-s}=|\alpha|^{-s}e^{-is(\Arg{\alpha})},-\pi<\Arg{\alpha}\le\pi
\end{equation*}
converges absolutely for $\mbox{Re}{s}\gg0$ and has an analytic continuation to the half plane $\mbox{Re}{s}>-\epsilon$ for some $\epsilon>0$ which is holomorphic at $s=0$.

The following theorem implies that the regularized determinants $\det_{\infty}(s \cdot \mathrm{id} - \Theta | H^{i}_{\Fh}(M,\C))$ satisfy the condition (1) for $s\in\C$. 

\begin{theorem}
\label{thm:3.1}
$\Sp(\Theta|H^{i}_{\Fh}(M,\C))$ can be described in terms of $\Sp(\varphi^{*}|H^{i}(S,\C))$
\begin{equation}
\label{eq_eigenvalues}
    \Sp(\Theta|H^{i}_{\Fh}(M,\C))=\left\{ \frac{\log \alpha + 2{\pi}{i}{v}}{\log{r}} |\alpha\in\Sp(\varphi^{*}|H^{i}(S,\C)), v\in\Z \right\}.
\end{equation}

\end{theorem}

\begin{proof}
Let $[\omega_{\Fh}] \in H^{i}_{\Fh}(M,\C)$ be an eigenvector of the infinitesimal operator $\Theta$ with an eigenvalue $\alpha\in\C$:
\begin{equation*}
     \Theta [\omega_{\Fh}] =\alpha\cdot[\omega_{\Fh}].
\end{equation*}
By using Theorem \ref{thm:2.1}, we have the following differential equation:
\begin{equation*}
    \frac{d}{dt}\Psi([\omega_{\Fh}])(t)=\alpha\cdot\Psi([\omega_{\Fh}])(t)
\end{equation*}
Then the solution is given as a cohomological class on $S$.
\begin{equation*}
    \Psi([\omega_{\Fh}])(t)=e^{\alpha{t}}\Psi([\omega_{\Fh}])(0)
\end{equation*}
If we note that $\Psi([\omega_{\Fh}])(t+\log{r})=\varphi^{*} (\Psi([\omega_{\Fh}])(t))$, $\Psi([\omega_{\Fh}])(t)$ is an eigenvector of $\varphi^{*}$ with the eigenvalue $e^{\alpha\log{r}}=r^{\alpha}$ for a fixed $t$ in $S^{1}$:
\begin{equation*}
    \begin{split}
    \Psi([\omega_{\Fh}])(t+\log{r})&=e^{\alpha\log{r}} \Psi([\omega_{\Fh}])(t)\\
    &=\varphi^{*} (\Psi([\omega_{\Fh}])(t)).
    \end{split}
\end{equation*}
Thus, for an eigenvalue $\alpha \in \Sp(\Theta|H^{i}_{\Fh}(M,\C))$ , $r^{\alpha}=e^{\alpha \log{r}}$ should be the eigenvalue of $\varphi^{*}$ on the $H^{i}(S,\C)$.

Conversely, for an eigenvector $[\mu] \in H^{i}(S,\C)$  of $\varphi^{*}$ with an eigenvalue $\beta$, we suspend it to foliation cohomological classes $[\mu_{\Fh}]_{v}$ for all $v\in\Z$ which are defined by:
\begin{equation*}
    \Psi([\mu_{\Fh}]_{v})(t) := e^{(\frac{\log\beta+2\pi {i}{v}}{\log{r}})t}[\mu]\ \mathrm{for} \ t \in \R/\log{r}\,\Z,
\end{equation*}
where $i$ is $\sqrt{-1}$. They satisfy the boundary condition and are eigenvectors of $\Theta$ on $H^{i}_{\Fh}(M,\C)$ with the eigenvalues $\frac{\log \beta + 2 \pi{i}{v}}{\log{r}}$ for $\forall v \in \Z$. Therefore, $\{ \frac{\log \beta + 2 \pi{i}{v}}{\log{r}} |v \in \Z \}$ is the corresponding eigenvalues of the eigenvalue $\beta$.
\end{proof}
To show that the condition (2) holds, we interpret the Dirichlet series by using the Hurwitz zeta function 
\begin{equation}\label{eq:3.2.1}
\begin{split}
    \tr(s \cdot \mathrm{id} - \Theta | H^{i}_{\Fh}(M,\C))^{-z}&=\sum_{\alpha\in\Sp(\varphi^{*}_{i})}\sum_{v\in\Z}(s-\frac{\log{\alpha}+2\pi{i}{v}}{\log{r}})^{-z}\\
    &=\sum_{\alpha\in\Sp(\varphi^{*}_{i})}\sum_{v\in\Z} \left\{ \frac{2\pi{i}{v}}{\log{r}}\left(\frac{s\cdot\log{r}-\log \alpha}{2{\pi}{i}}+v\right) \right\}^{-z},
\end{split}
\end{equation}
where $\varphi^{*}_{i}$ denotes $\varphi^{*}$ on $H^{i}(S)$.

We recall that the Hurwitz zeta function $\zeta_{\mathfrak{hur}}(z,s)$ is  defined by:
\begin{equation*}
    \zeta_{\mathfrak{hur}}(z,s):=\sum_{n=0}^{\infty}(s+n)^{-z}\,\mathrm{for}\,\mathrm{Re}(z)>1,\,\mathrm{Re}(s)>0.
\end{equation*}
It is known that the series is absolutely convergent and can be extended to a meromorphic function defined for all $z\ne1$.
More generally we set
\begin{equation*}
    \zeta_{\eta,\mathfrak{hur}}(z,s)=\sum_{n=0}^{\infty}(\eta(s+n))^{-z},
\end{equation*}
where $\eta$ is a non-zero complex number.
If we denote  $\frac{2 \pi i}{\log{r}}$ and $\frac{s\cdot\log{r}-\log \alpha}{2{\pi}{i}}$ in (\ref{eq:3.2.1}) simply by $\eta$ and $s_{\alpha}$, respectively, then the series (\ref{eq:3.2.1}) is expressed by using the Hurwitz zeta function
\begin{equation*}
    \sum_{\alpha\in\Sp(\varphi^{*}_{i})}\zeta_{\eta,\mathfrak{hur}}(z,s_{\alpha})+\zeta_{-\eta,\mathfrak{hur}}(z,-s_{\alpha})-(\eta{s_{\alpha}})^{-z}.
\end{equation*}
Since the condition (2) holds from the property of the Hurwitz zeta function, the regularized determinant is defined for $(M,\Fh_{M},\phi^{t})$.


Next we show that the leafwise cohomological expression coincides with the dynamical zeta function of $(M,\Fh_{M},\phi^{t})$. We define a regularized product by 
\begin{equation*}
    \begin{split}
    \prod_{v\in\Z}\eta(s+v)&=\prod_{n\in\N}\eta(s+n)\cdot(-\eta)(-s+n)/(\eta{s})\\
    &:=\exp(-\partial_{z}\zeta_{\eta,\mathfrak{hur}}(0,s)-\partial_{z}\zeta_{-\eta,\mathfrak{hur}}(0,-s))\cdot(\eta{s})^{-1},
    \end{split}
\end{equation*}
where $\eta$ is a non-zero complex number. Then the lemma follows from properties of Hurwitz zeta function.
\begin{lemma}
\label{lem:3.2}
The regularized product such as the following form can be computed by using properties of the Hurwitz zeta functions:
\begin{equation*}
    \prod_{v\in\Z}\eta(s+v)=1-e^{-2\pi{i}s}
\end{equation*}
\end{lemma}
\begin{proof}
It is known for the Hurwitz zeta functions that:
\begin{equation*}
    \partial_{z}\zeta_{\mathfrak{hur}}(0,s)=\log\Gamma(s)-\frac{1}{2}\log2\pi.
\end{equation*}
Then we have
\begin{equation*}
    \exp(-\partial_{z}\zeta_{\eta,\mathfrak{hur}}(0,s))=\eta^{\frac{1}{2}-s}\left( \frac{\Gamma(s)}{\sqrt{2\pi}} \right)^{-1}.
\end{equation*}
By using the formula:
\begin{equation*}
    \frac{1}{s}\left( \frac{\Gamma(s)}{\sqrt{2\pi}} \right)^{-1}\left( \frac{\Gamma(-s)}{\sqrt{2\pi}} \right)^{-1}=e^{\pi{i}(\frac{1}{2}+s)}(1-e^{-2\pi{i}{s}}),
\end{equation*}
we get the following equality:
\begin{equation*}
    \begin{split}
        \prod_{v\in\Z}\eta(s+v)&=\eta^{\frac{1}{2}-s}\left( \frac{\Gamma(s)}{\sqrt{2\pi}} \right)^{-1}\cdot(-\eta)^{\frac{1}{2}+s}\left( \frac{\Gamma(-s)}{\sqrt{2\pi}} \right)^{-1}\cdot(\eta{s})^{-1}\\
        &=(\eta)^{-\frac{1}{2}+s}(-\eta)^{\frac{1}{2}+s}\cdot\frac{1}{s}\left( \frac{\Gamma(s)}{\sqrt{2\pi}} \right)^{-1}\left( \frac{\Gamma(-s)}{\sqrt{2\pi}} \right)^{-1}\\
        &=\begin{cases}
        1-e^{-2\pi{i}{s}}&\,\mathrm{If}\,\,0\leq\Arg{\eta}<\pi\\
        1-e^{2\pi{i}{s}}&\,\mathrm{If}\,\,-\pi\leq\Arg{\eta}<0
        \end{cases}
    \end{split}
\end{equation*}
In this case, since $0\leq\Arg{\eta}<\pi$, we get the statement of the lemma.
\end{proof}
We denote the leafwise cohomological expression by the regularized product as follows
\begin{equation*}
\begin{split}
    \prod_{i}^{d}\det_{\infty}(s\cdot\mbox{id}-\Theta|H^{i}_{\Fh}(M,\C))^{(-1)^{i+1}}&\overset{(\ref{eq_eigenvalues})}{=}\prod_{i}\prod_{\alpha\in\Sp(\varphi^{*}_{i})}\prod_{v\in\Z} (s-\frac{\log{\alpha}+2\pi{i}{v}}{\log{r}})^{(-1)^{i+1}}\\
    &=\prod_{i}\prod_{\alpha\in\Sp(\varphi^{*}_{i})}\prod_{v\in\Z}  \left\{ \frac{2\pi{i}{v}}{\log{r}}\left(\frac{s\cdot\log{r}-\log \alpha}{2{\pi}{i}}+v\right) \right\}^{(-1)^{i+1}}.
\end{split}
\end{equation*}
By using the Lemma \ref{lem:3.2}, we have the following product form
\begin{equation*}
    \prod_{i}^{d}\det_{\infty}(s\cdot\mbox{id}-\Theta|H^{i}_{\Fh}(M,\C))^{(-1)^{i+1}}=\prod_{i}\prod_{\alpha\in\Sp(\varphi^{*}_{i})}(1-\exp(\log\alpha-s\cdot\log{r}))^{(-1)^{i+1}}.
\end{equation*}
The corollary follows from the product form.
\begin{corollary}
\label{cor:3.2.1}
The leafwise cohomological expression is described with the cohomology of the fiber $S$:
\begin{equation}\label{eq:3.5}
    \prod_{i}^{d}\det_{\infty}(s\cdot\mbox{id}-\Theta|H^{i}_{\Fh}(M,\C))^{(-1)^{i+1}} =\prod_{i=0}^{d}\det(1-\varphi^{*}\cdot r^{-s}|H^{i}(S,\C))^{(-1)^{i+1}}.
\end{equation}
\end{corollary}

The relation between periodic orbits of $\Z$-action on $S$ and $\R$-action on $M$ gives the following lemma.

\begin{lemma}
the dynamical zeta function $\zeta(M;s)$ of $(M,\Fh_{M},\phi^{t})$ has the cohomological expression in terms of the cohomology groups of the fiber $S$
\begin{equation}\label{eq:3.6}
    \zeta(M;s) =\prod_{i=0}^{d}\det(1-\varphi^{*}\cdot r^{-s}|H^{i}(S,\C))^{(-1)^{i+1}}.
\end{equation}
\end{lemma}

\begin{proof}
We have the bijective between periodic orbits of $\Z$-action on $S$ and $\R$-action on $M$.
If we denote by $\mathfrak{o}$ a corresponding periodic orbit of $\Z$-action on $S$ for a periodic orbit $\gamma$ of $\R$-action on $M$, the norm $N(\gamma)$ is defined by $\log{N(\gamma)}=|\mathfrak{o}|\log{r}$ where $|\mathfrak{o}|$ is the period of $\mathfrak{o}$.
Therefore we have
\begin{equation*}
    \zeta(M;s)=\prod_{\mathfrak{o}}(1-e^{-s|\mathfrak{o}|\log{r}})^{-1},
\end{equation*}
where $\mathfrak{o}$ runs over periodic orbits of the $\Z$-action on $S$.
Then the lemma follows from the Lefschetz fixed-point theorem.
\end{proof}

\begin{proof}[Proof of theorem 3.1]
The assertion follows from the composite of (\ref{eq:3.5}) and (\ref{eq:3.6}).
\end{proof}

\subsection{Functional equations}

The zeta function $\zeta(M;s)$ has the functional equation as follows:
The closed manifold $S$ has the symmetric property which is given by a perfect pairing
\begin{equation*}
    \cup:H^{i}(S) \times H^{d-i}(S) \rightarrow H^{d}(S).
\end{equation*}
The pullback of the diffeomorphism $\varphi$ of $S$ is compatible with the pairing
\begin{equation*}
    \varphi^{*}(v)\cup\varphi^{*}(w)=\varphi^{*}(v\cup{w})\quad\mathrm{for}\,v\in{H}^{i}(S),w\in{H}^{d-i}(S).
\end{equation*}
Since $\varphi^{*}$ acts on $H^{d}(S)$ as the identity, the following holds (\cite{Ha},Appendix C, Lemma 4.3.):
\begin{equation}\label{eq:3.3}
    \det(\varphi^{*}|H^{d-i}(S))=\frac{1}{\det(\varphi^{*}|H^{i}(S))}
\end{equation}
and
\begin{equation}\label{eq:3.4}
    \det(1-\varphi^{*}t|H^{d-i}(S))=\frac{(-t)^{\beta(i)}}{\det(\varphi^{*}|H^{i}(S))}\det(1-\varphi^{*}t^{-1}|H^{i}(S)),
\end{equation}
where $\beta(i)$ is the $i$-th Betti number.
\begin{theorem}[Functional equation]
\label{functional_eq}
The zeta function is symmetrical about $Re(s)=0$
\begin{equation*}
    \zeta(M;s)=(-r^{s})^{\chi(S)} \cdot \zeta(M;-s),
\end{equation*}
where $\chi(S)$ is the Euler characteristic of $S$.
\end{theorem}
\begin{proof}
As the result of the corollary \ref{cor:3.2.1}, the zeta function can be described with the cohomology groups of $S$:
\begin{equation*}
    \begin{split}
            \zeta(M;s)&=\prod_{i=0}^{d}\det(1-\varphi^{*}\cdot r^{-s}|H^{d-i}(S,\C))^{(-1)^{i+1}}\\
            &\overset{(\ref{eq:3.4})}{=}\prod_{i=0}^{d}\left\{\frac{(-r^{-s})^{\beta(i)}}{\det(\varphi^{*}|H^{i}(S))}\det(1-\varphi^{*}r^{s}|H^{i}(S))\right \}^{(-1)^{i+1}}\\
            &\overset{(\ref{eq:3.3})}{=}(-r^{-s})^{-\chi(S)}\zeta(M;-s).
    \end{split}
\end{equation*}
\end{proof}



\subsection{Special values of \texorpdfstring{$\zeta(M;s)$}{}}

We compute special values of the dynamical zeta function $\zeta(M;s)$:
We define the special value of $\zeta(M;s)$ at $s=k$ by
\begin{equation*}
    \zeta(M;k)^{*}:=\lim_{s\rightarrow k}\zeta(M;s)\cdot (s-k)^{-\mathrm{ord}_{s=k}\zeta(M;s)},
\end{equation*}
where ${-\mathrm{ord}_{s=k}\zeta(M;s)}$ is the order of $\zeta(M;s)$ at $s=k$ for $k\in\C$.
Note that the order of $\zeta(M;s)$ at $s=k$ for $k\in\C$ is given by
\begin{equation*}
    \mathrm{ord}_{s=k}\zeta(M;s)=\sum_{i}(-1)^{i+1}\dim(H^{i}_{\Fh}(M)^{\Theta\sim k}),
\end{equation*}
where $H^{\sigma\sim{k}}$ is the eigenspace of $\sigma$ on $H$ whose eigenvalue is $k$. 


We show the computation of special values of $\zeta(M;s)$ as follows.
\begin{theorem}
If we set the $m$-th Lefschetz number for $m\in\Z$ by 
\begin{equation*}
    \Lambda(\varphi^{m}):=\sum_{i}(-1)^{i}\mathrm{tr}(\varphi^{m*}|H^{i}(S)),
\end{equation*}        
the special value of $\zeta(M;s)$ at $s=k$ can be expressed  with the order of the function and the Lefschetz number.
\begin{equation*}
    \zeta(M;k)^{*}=(\log{r})^{\mathrm{ord}_{s=k}\zeta(M;s)}\exp(\sum_{m\ge1}\frac{r^{-km}\Lambda(\varphi^{m})+\mathrm{ord}_{s=k}\zeta(M;s)}{m}).
\end{equation*}
\end{theorem}
\begin{proof}
Corollary \ref{cor:3.2.1} allows $\zeta(M;s)$ to be decomposed with respect to eigenvalues of $\varphi^{*}$:
\begin{equation*}
    \zeta(M;s)=\prod_{i}(1-r^{k-s})^{(-1)^{i+1}\dim (H^{i}(S))^{\varphi^{*}\sim r^{k}}}\det(1-\varphi^{*}\cdot{r}^{-s}|H^{i}(S)^{\varphi^{*}\nsim r^{k}})^{(-1)^{i+1}}
\end{equation*}
We give the Taylor series about $(s-k)$
\begin{equation*}
    \begin{split}
    \zeta(M;s)=(\log{r}(s-k)+&O((s-k)^{2}))^{\mathrm{ord}_{s=k}\zeta(M;s)}\\
    &\cdot\exp(\sum_{m\ge1}\sum_{i}(-1)^{i}\mathrm{tr}(\varphi^{m*})\cdot \frac{r^{-sm}}{m}|H^{i}(S)^{\varphi^{*}\nsim r^{k}}).
    \end{split}
\end{equation*}
Note that $O(x^{2})$ is the big-O notation. Then we compute the special value at $s=k$
\begin{equation*}
    \zeta(M;k)^{*}=(\log{r})^{\mathrm{ord}_{s=k}\zeta(M;s)}\cdot\exp(\sum_{m\ge1}\sum_{i}(-1)^{i}\mathrm{tr}(\varphi^{m*})\cdot \frac{r^{-km}}{m}|H^{i}(S)^{\varphi^{*}\nsim r^{k}}).
\end{equation*}
By decomposing the Lefschetz number with respect to $\mathrm{ord}_{s=k}\zeta(M;s)$
\begin{equation*}
        \Lambda(\varphi^{m})=-r^{km}\cdot\mathrm{ord}_{s=k}\zeta(M;s)+\sum_{i}(-1)^{i}\mathrm{tr}(\varphi^{m*}|H^{i}(S)^{\varphi^{*}\nsim r^{k}}),
\end{equation*}
we get the statement of the theorem
\begin{equation*}
    \zeta(M;k)^{*}=(\log{r})^{\mathrm{ord}_{s=k}\zeta(M;s)}\exp(\sum_{m\ge1}\frac{r^{-km}\Lambda(\varphi^{m})+\mathrm{ord}_{s=k}\zeta(M;s)}{m}).
\end{equation*}
\end{proof}